\section{Conclusion}
In this work, we explored the utility of register token embeddings as auxilliary features that supplement the \texttt{[CLS]} token in Vision Transformers (ViTs) to improve the feature richness for out-of-distribution (OOD) generalization and anomaly rejection. While previous approaches discard register tokens after training, we demonstrated that these tokens contain valuable information that augments the global features captured by the \texttt{[CLS]} token. By concatenating the \texttt{[CLS]} token with the mean of register tokens, we produced richer representations for the linear classifier to perform better on OOD tasks across multiple ViT architectures pre-trained with or without registers. Our approach consistently outperformed baseline methods in both generalization and anomaly rejection, achieving notable improvements in top-1 accuracy and reductions in false positive rates. Importantly, these gains were achieved without introducing additional computational overhead. These findings suggest that register tokens are not redundant but instead play a crucial role in enhancing the robustness and adaptability of ViTs. 
